\section{Objetivos del curso}

\begin{frame}
	\frametitle{\insertsection}

	\begin{columns}
		\begin{column}{.45\paperwidth}
			Al final de este curso, ustedes serán capaces de

			\begin{enumerate}
				\item \textbf{Comprender} la teoría detrás de los métodos numéricos.
				      \begin{itemize}
					      \item Consistencia, convergencia, estabilidad.
					      \item Análisis de errores.
				      \end{itemize}

				\item \textbf{Implementar} esquemas numéricos en Python.
				      \begin{itemize}
					      \item Métodos de diferencias finitas (FDM).
					      \item Métodos de elementos finitos (FEM).
					      \item Métodos de volúmenes finitos (FVM).
				      \end{itemize}

				\item \textbf{Resolver} problemas prácticos.
				      \begin{itemize}
					      \item Ecuaciones parabólicas (calor).
					      \item Ecuaciones elípticas (Poisson).
					      \item Ecuaciones hiperbólicas (ondas, leyes de conservación).
				      \end{itemize}
			\end{enumerate}
		\end{column}
		\begin{column}{.4\paperwidth}
			\begin{figure}[ht!]
				\centering
				\animategraphics[loop,controls,height=.5\paperheight]{10}{hellopetsc-}{0}{99}
				\caption{Simulación numérica de la ecuación de calor unidimensional.}
				\label{fig:calor}
			\end{figure}
		\end{column}
	\end{columns}
\end{frame}
